% ─────────────────────────────────────────────────────────────────────
% Trajectory-level escalation: design (input after \S routing) and
% evaluation (input before \S discussion). Numbers come from
% benchmarks/reports/harness-live-mbpp-main-*.json (see benchmarks/README.md).
% ─────────────────────────────────────────────────────────────────────

\subsection{Trajectory-Level Escalation}
\label{sec:escalation}

The router of \S\ref{sec:routing} decides \emph{before} work starts, from a self-reported complexity estimate.
Cascades~\cite{chen2023frugalgpt} and learned routers~\cite{ong2024routellm,ding2024hybrid} make the same decision at the granularity of a whole request.
For multi-step agent tasks this granularity is wrong: a small model can usually execute most of a trajectory and fails on one step, and the cost of a wrong decision is paid over the entire trajectory.
Joule therefore makes the small model the \emph{executor} and decides after every step whether to keep going, ask for help, or hand over.

\paragraph{Execution state.}
A run maintains a structured state $S$: the goal, a versioned plan, completed steps with their tool results, observations, failures keyed by a normalized error signature, hypotheses, advice received, and every escalation decision.
$S$ is the only input to escalation; the large model never restarts from the original prompt.

\paragraph{Step agent.}
One model-driven loop serves every tier.
At each turn the current model proposes exactly one structured action---a tool call with an optional verifier, a final answer, a focused question, or a give-up---and the runtime executes it, verifies it, and updates $S$.

\paragraph{Deterministic verification.}
Wherever possible a step is verified without a model: an output pattern, a DOM predicate, or a command exit code with a parsed pass fraction (``3/5 tests passed'').
A command that exits non-zero fails verification even when the tool call itself succeeded; a search command that exits 1 to report ``no match'' does not.
A source file the agent has just written is compile-checked before anything else happens; a syntax error goes back to the agent with the offending line and counts as a slip rather than as being stuck, escalating only if the same error repeats.
Model-judged verification is opt-in and always labelled, so results can be reported separately.

\paragraph{Evidence-based confidence.}
No model self-report enters the confidence used for escalation.
After each step,
\begin{equation}
c = w_t\,\mathrm{tool} + w_v\,\mathrm{verify} + w_p\,\mathrm{progress} + w_b\,\mathrm{budget} + w_c(1-\mathrm{repeat}) - w_r\,\mathrm{repeat} - w_x\,\mathrm{contradiction},
\end{equation}
where $\mathrm{tool}$ is the last tool outcome, $\mathrm{verify}$ the last verifier outcome, $\mathrm{progress}$ the verified share of a recent window (an improving pass fraction counts), $\mathrm{repeat}$ the recurrence of a failure signature, and $\mathrm{contradiction}$ indicates a tool success whose verification failed.

\paragraph{Escalation policy.}
An ordered rule set maps $(S, c)$ to one of four actions:
\textsc{continue} (same tier keeps working);
\textsc{consult} (one focused question to the large model, whose answer returns to the small model as advice);
\textsc{handoff} (the large model continues from a context rendered from $S$: goal, plan, completed work, key observations, failures, advice, and unresolved questions);
\textsc{abort} (budget, safety, or an impossible tool requirement).
Hard triggers precede soft thresholds: repeated malformed output, a give-up, or $k$ accumulated failures force a handoff; the same failure signature twice, a stall, or a verification failure that survives a retry without improvement triggers a consult; thresholds on $c$ apply only once two failures have accumulated.
Every escalation is gated by the budget envelope (\S\ref{sec:budget}): a consult must fit in the remaining cost and tokens, and a handoff also consumes an escalation unit.

\paragraph{Ladder.}
Tiers form a ladder $\langle \text{slm}, \text{mid}, \text{llm} \rangle$ in which the optional middle rung is an efficient large model.
\textsc{consult} asks the next rung up and \textsc{handoff} moves execution to it, including after a reasoning breakdown (repeated unparseable output, or a give-up before any step succeeded); an option sends breakdowns straight to the top rung, and \S\ref{sec:eval-escalation} shows why the default climbs one rung instead.
After a handoff the policy is rung-local: the new model is judged on its own steps, failures, and consultations, not on the record of the model it replaced.
Failures are counted within a window of recent steps, so a long task that keeps making verified progress is not mistaken for a stuck one.
Rungs no provider serves are dropped, so a two-model deployment degrades to the two-tier policy.

\paragraph{Why consult.}
A handoff bills the remaining trajectory at large-model prices.
A consult bills one question and one answer, and returns control to the small model.
When the small model is stuck on a single decision, which is the common case on coding tasks, a consult recovers the trajectory at a fraction of a handoff's cost.
We measure how often it does.

\paragraph{Consultations that return the fix.}
A consultation originally returned prose.
On long tasks the small model turned correct advice into an incorrect edit often enough that most consultations ended in a handoff anyway.
The consultant therefore sees the current contents of the files the run has touched and may return, alongside its answer, a concrete edit: a whole-file replacement for small files or one exact search-and-replace for large ones.
The runtime applies the edit through the same write tool the agent uses, compile-checks it, records it as a step attributed to the advisor's tier, and returns control to the small model, which verifies and finishes.
The consultation remains a single call; what changes is that the small model no longer has to transcribe the fix.

\paragraph{Switches for ablation.}
Each design choice above maps to one policy switch, so the benchmark harness can remove exactly that choice: no consultations (every escalation is a handoff), prose-only consultations, no step verification, the model's self-reported confidence in place of the evidence-based composite, and no static checks.

% ─────────────────────────────────────────────────────────────────────
